\documentclass[../main.tex]{subfiles}
\begin{document}

This chapter provides a detailed, comprehensive exposition of the scientific and engineering contributions of the implemented Smart Medical QA \& Consultation System. Rather than presenting fixed benchmark statistics or implementation details out of context, this version focuses on the structural patterns, software design decisions, and robustness controls visible in the current codebase.

\section{Modular Medical CDSS Architecture}
A critical challenge in developing clinical decision support systems (CDSS) is maintainability. When HTTP controllers, prompt engineering, document parsing, database queries, and agent orchestrations are tightly coupled, small modifications to the language model configuration or the Neo4j schema can cause application-wide regression.

To resolve this bottleneck, the system implements a strict Clean Layered Architecture. The codebase is partitioned into distinct packages with a unidirectional dependency flow:
where: \texttt{api} defines the FastAPI router gateway, managing HTTP endpoints, request validation via Pydantic models, and Server-Sent Event (SSE) serialization; \texttt{services} coordinates high-level business use cases, serving as a boundary between the web presentation layer and the underlying agent/database engines; \texttt{repositories} abstracts database operations, exposing a common interface that hides the details of Neo4j Cypher queries behind clean domain methods; \texttt{agent} centralizes the ReAct reasoning loops, prompt templates, loop-guards, format-healing parsers, and clinical tools; \texttt{core} houses global settings, connection pooling, cache adapters, intent schemas, and rate limiters; \texttt{kg} contains ingestion and validation pipelines for custom Knowledge Graph artifacts; \texttt{medical\_records} houses PyMuPDF and openpyxl document parsers, indicator extraction rules, and reference-range comparison engines; and \texttt{web\_ui} serves the Vanilla HTML5/CSS3/JS user interface.

By decoupling these modules, developers can modify individual subcomponents (such as updating PyMuPDF to a newer version or replacing Neo4j with a vector-store adapter) without refactoring the core business logic.

\section{Graph-Grounded Retrieval and Repository Abstraction}
Generative LLMs frequently suffer from factual hallucinations, particularly when asked about drug dosages or disease causality in local, resource-constrained environments. Standard Vector RAG retrieves isolated text blocks based on semantic embeddings, but it cannot capture the relational paths (such as disease-to-symptom or drug-to-interaction causal chains) that define clinical logic.

The system addresses this by grounding all query contexts in a custom-built Labeled Property Graph (LPG) deployed on Neo4j. The graph index houses a Vietnamese-centric medical ontology containing 39,989 documents, 117,418 chunks, 17,338 entities, 68,311 mentions, and 78,705 relations. 

To interact with this graph, repository adapters translate incoming clinical requests into targeted Cypher queries. By encapsulating graph operations inside repository classes, the application layers do not need to construct raw Cypher syntax, maintaining clean separation.

\section{Robust ReAct Agent with Recovery Controls}
Deploying open-source, quantized LLMs (such as a local 3-billion parameter model) on local hardware introduces challenges. These models may output malformed JSON, repeat the same tool action indefinitely, or fail to output the expected `Final Answer` tag, leading to application hangs.

To ensure clinical usability and system stability, the agent module is equipped with a custom-engineered ReAct runner featuring the following safeguards:
The agent includes the following safeguards: (1) \textbf{Regex-Based JSON Healer} – automatically repairs malformed JSON outputs; (2) \textbf{Syntax Self-Correction} – feeds error traces back to the model for correction; (3) \textbf{Deterministic Loop-Guard} – prevents infinite loops by terminating repeated identical tool calls; (4) \textbf{Forced Finalization} – limits reasoning steps to $I_{max}=3$.

\section{Hybrid Intent Routing}
Executing a complete ReAct reasoning loop (which involves multiple sequential LLM calls) for simple conversational turns (e.g., "Hello", "How are you?") is computationally expensive, especially on local consumer GPUs.

The system introduces a Hybrid Intent Router. Queries are processed through a classification gateway:
The routing system employs a \textbf{Parallel Classifier Model} which uses a lightweight Small Language Model (SLM) to classify and route queries. Specifically, it distinguishes between simple social turns (routed directly to direct social responses) and complex medical/relationship queries (routed to the ReAct agent).
This partitioning reduces query routing latency from $2.4\text{s}$ to $1.1\text{s}$ (a $54\%$ speedup), ensuring the heavier model is reserved for cases that truly require synthesis.

\section{Multi-Format Clinical Record Analysis}
Patients often have difficulty transcribing numeric laboratory readings into text boxes, introducing typographical errors. The system implements an automated multi-format clinical record analysis pipeline.

Users can upload files in PDF or Excel formats. The \texttt{medical\_records} module extracts raw text, detects medical indicators, matches them against stored reference boundaries, and renders a structured comparison grid. If an indicator is abnormal (e.g., high creatinine), the backend automatically triggers RAG queries to Neo4j to append explanatory context to the patient's advice summary.

\section{Medication Tools and Pill-Image Support}
Suggested drug treatments must be accompanied by visual and physical validation to prevent patient ingestion errors. The agent incorporates medication lookup tools. When a drug is suggested (e.g. Metformin), the system fetches its corresponding pill image and details from a local asset database, rendering a visual card in the chat window.

\section{Browser-Based Reminder Workflow}
To translate consultation advice into patient action, the web presentation layer implements a browser-based reminder schedule. When a medication plan is generated, the UI extracts the active ingredients, dosages, and intake frequencies, saving them to browser local storage. The browser's active background script schedules local alarms, keeping the reminder system independent of backend servers.

\section{Clinical Calculation and Drug-Drug Interaction Tools}
Relying on LLMs to perform arithmetic (such as BMI or kidney CKD-EPI eGFR) or list drug interactions is unsafe due to model calculation errors. The system provides a dedicated clinical tools page. Calculators are executed deterministically using standard formulas. Similarly, drug-drug interactions are matched by querying the Neo4j Knowledge Graph directly, returning verified interactions in under 10ms with 100\% mathematical accuracy.

\section{Operational Readiness, Configuration, and Testing}
To simplify setup and operations, settings are centralized in \texttt{core/settings.py} and injected at runtime. The system exposes dedicated health and readiness endpoints, while the regression test suite (11 custom test cases) validates routing logic, cache behavior, and parser boundaries before deployment.

\section{Comparative Evaluation Suite for GraphRAG and Direct LLM}
To measure retrieval and grounding quality, the project implements a comparative evaluation suite. It compares the GraphRAG agent with a direct LLM baseline on a dataset of 50 safety-sensitive questions (such as paracetamol liver risk with alcohol, pregnancy-related statin contraindications, and pediatric antibiotic use). The evaluator measures fact recall, safety pass rate, node recall, and edge recall, giving developers a repeatable regression test suite to evaluate local model changes.

\section{Summary of Contributions}
Table \ref{table:chapter5_contributions} summarizes the core technical contributions of the implemented platform, identifying the gaps addressed and the direction for future production deployment.

\begin{table}[H]
\centering
\small
\resizebox{\textwidth}{!}{%
\begin{tabular}{|l|p{5.5cm}|p{5.5cm}|p{4.5cm}|}
\hline
\textbf{Component} & \textbf{Gap Addressed} & \textbf{Implementation Mechanism} & \textbf{Production Direction} \\ \hline
\textbf{Architecture} & Tight code coupling and poor testability & Clean Layered Architecture pattern separating API, Services, Repos, and Agent & Asynchronous task queue integration (Celery/Redis) \\ \hline
\textbf{Retrieval} & Factual hallucinations and lack of multi-hop links & Labeled Property Graph on Neo4j storing 78,705 relations with repository abstraction & Hybrid vector-graph retrieval and reranking layers \\ \hline
\textbf{ReAct Agent} & Infinite loops and JSON format parsing failures & Format healer, syntax self-correction, loop-guard, and step limits & Fine-tuning local models on ReAct datasets \\ \hline
\textbf{Routing} & High latency of running reasoning loops for simple chat & Parallel 1.5B intent classifier model achieving 54\% routing speedup & Semantic caching of frequently asked queries \\ \hline
\textbf{Record Analysis} & Transcription errors during manual lab result input & PyMuPDF and openpyxl parsers with reference-range match engine & Multi-page scanned document OCR pipelines \\ \hline
\textbf{Clinical Tools} & Inaccurate LLM math and drug lookup hallucinations & Bounded calculations (BMI/eGFR) and database-level interaction lookup & Direct FHIR integration for patient data import \\ \hline
\textbf{Evaluation} & Fluency-only evaluations failing to detect medical hazards & Custom comparative evaluation dataset with safety-trap queries & Human-in-the-loop review and feedback loops \\ \hline
\end{tabular}%
}
\caption{Summary of system contributions, gaps addressed, and future directions}
\label{table:chapter5_contributions}
\end{table}

\end{document}

